\clearpage
\section{Conclusiones}
\label{sec:conclusions-final}
\begingroup
\fontsize{10.2}{12.8}\selectfont
\setlength{\parskip}{3pt plus .5pt}

La tesis geométrica queda establecida mediante reconstrucción y transporte:
una representación positiva de la estructura discreta del continuo conserva,
con sus marcas, el registro del que procede y los lectores que éste
determina. La carga total y los cocientes de menores recuperan cada
coordenada; las dos memorias nonádicas proporcionan otra reconstrucción
operatoria. El valor y la forma se articulan así por inversas explícitas,
después de la generación APP--TRIT--TPK y de sus prolongaciones compatibles.
La doble proyección mantiene unidas la realización aritmética y la
incidencia excepcional dentro de esa genealogía.

El refinamiento supera la limitación de una matriz con trece columnas.
Para cualquier rango y dimensión externa finitos, un nivel nonádico
suficiente produce los menores positivos y la imagen amplituhedral
correspondiente. Los nodos heredados conservan sus banderas; los
promedios celulares conservan la medida y registran el detalle
ortogonal. En rango uno, la cobertura simplicial y la recursión de
residuos se extienden a toda dimensión finita. En rangos mayores,
el resultado determina colecciones positivas compatibles con rango y
límite probados; su amplitud dimensional queda separada de la
clasificación global de todas las celdas amplituhedrales.

La conservación admite tres expresiones complementarias. La forma
canónica se suma bajo subdivisión de un soporte fijo; la energía
efectiva se conserva mediante reducción de Schur; y la medida
cronológica determina momentos cuyo incremento se descompone en
detalles recuperables. Estos momentos construyen operadores de
Jacobi cuyas compresiones finitas tienen espectro simple y residuos positivos. Se obtiene así
un recorrido efectivo desde las separaciones aritméticas hasta una
representación espectral, con fracción continua determinada por la
misma medida. Las mutaciones que cambian únicamente la coordenatización
conservan los lectores del punto marcado.

La realización física gana precisión mediante sus operadores. Los
Hamiltonianos de los tres regímenes tríticos poseen formas
simplécticas y transformaciones de Legendre explícitas. La coordenatización
elíptica variable incorpora su término de conexión; la memoria
parabólica conserva su invariante; la apertura hiperbólica tiene
su evolución propia. Los lectores de acción y respuesta
constitutiva se transportan por la inversa grassmanniana. Con
las representaciones regionales y las marcas enteras declaradas,
la respuesta constitutiva etiquetada recupera a su vez el registro.
La condición de igualdad de trazas gobierna el pegado de formas
con valores en la recta de acción.

La deformación excepcional conserva el volumen y la simetría
ternaria, y convierte la inversión del parámetro en dualidad
reticular. Su reciprocidad theta y su realización de signatura
partida enlazan la variación geométrica con la dualidad de la
red. En Yang--Mills, la escala areal--volumétrica, la medida común,
la coercividad uniforme y el testigo persistente sitúan la brecha
junto al Hamiltoniano que la realiza. El desarrollo conserva
el espacio físico y los entrelazadores necesarios para transportar
esa afirmación a través del refinamiento.

El resultado conjunto es una arquitectura de formas con memoria:
la dimensión puede aumentar, la representación puede cambiar y
la información pertinente permanece recuperable mediante mapas
determinados. Éste es el contenido diferencial del artículo.
La geometría positiva, la dualidad reticular y la realización
espectral quedan enlazadas por construcciones que especifican
qué generan, qué conservan y cómo se reconstruyen desde un mismo
origen discreto.
La representación de incidencia conserva también la componente uniforme
del registro: su inversa reconstruye las doce coordenadas y permite
recuperar la dirección geométrica después del transporte de memoria.
En la realización sectorial, el área total y la ocupación ponderada
determinan conjuntamente el reparto héxada--óctada. La purificación
especificada convierte ese reparto en una estructura de entrelazamiento,
con identidad finita entre entropía, áreas sectoriales y entropía relativa.
La conservación de información queda así precisada para cada aplicación:
registro completo, dirección, área y composición poseen inversas o
fibras explícitamente distintas.
La escala de área se determina mediante una longitud generada por la composición
lineal del retorno registrado y la norma local de Hadamard. La
reciprocidad radial determina el coeficiente gravitatorio en la
realización reducida declarada y reconoce después esa longitud como
longitud de Planck. Las mutaciones del mismo punto marcado y los
refinamientos heredados conservan el operador de área, sus
multiplicidades y la identidad modular. Las formas canónicas con
valores en ese operador extienden la conservación a los residuos,
mientras el radio de cada horizonte mantiene su factor geométrico
propio. Esta composición hace explícita la dependencia entre
longitud, incidencia y área sin alterar los estados sectoriales.
\par\endgroup
\clearpage
